\section{Terminal occupation and the localized upper bound}
\label{sec:terminal-upper}

The radial traversal estimates leave the occupation accumulated inside each terminal disk to be controlled. We first compare the entrance--exit chain in a centered block, then derive its exponential occupation moment. The last subsection combines that moment with the traversal barrier to prove the localized upper estimate.

\subsection{Endpoint products in one centered block}
\begin{lemma}\label{lem:typed-centered-endpoint-product}
Let \(D=D(y,eM)\), let \(A\) be the inner vertex boundary of
\(D(y,M)\), and let \(E\) be the directed edges leaving \(D\).  Put
\[
 H^S(a,e)=P_a((S_{\tau_D-1},S_{\tau_D})=e),\qquad
 ({\cal H}^Sf)(a)=\sum_eH^S(a,e)f(a,e).
\]
Let \(G^S=G_D\).  For \(x,x'\in V:=D(y,M/2)\), define on
\(A\times E\)
\[
 m_x(a,e)={G^S(a,x)H^S(x,e)\over H^S(a,e)},
\]
\[
\begin{aligned}
q_{x,x'}(a,e)=\frac{1}{H^S(a,e)}\big(&G^S(a,x)G^S(x,x')H^S(x',e)\\
  &+G^S(a,x')G^S(x',x)H^S(x,e)\\
  &-\mathbf1_{\{x=x'\}}G^S(a,x)H^S(x,e)\big),
\end{aligned}
\]
and let
\begin{equation}
 {\cal F}_M^{\rm G}=\{m_x/C:x\in V\}\cup
 \{q_{x,x'}/[C(1+G^S(x,x'))]:x,x'\in V\},               \tag{TP0}\label{eq:tu-TP0}
\end{equation}
where the fixed \(C\) is large enough that all functions have sup norm
at most one.
Let \(R_i^S(e,da')\) be a fixed radial-only centered return word from a
previous exit edge to the next inward entrance, let
\(\sigma_i=R_i^S1\), \(J_i^S=R_i^S/\sigma_i\), and set
\begin{equation}
 P_i^S((a,e),d(a',e'))=J_i^S(e,da')H^S(a',de').           \tag{TP1}\label{eq:tu-TP1}
\end{equation}
The word may be selected adaptively from its radial prefix, but not from an
angular endpoint.  If it contains \(q_i\) elementary letters, all at scales
at least \(M\), and \(\lambda_A^S\) is the center-to-\(A\) first-hit law,
put \(\Lambda^S=\lambda_A^SH^S\).  For \(f\in{\cal F}_M^{\rm G}\) and
every consecutive portion \(r\le s\) of
one centered block,
\begin{equation}
 \left|\Lambda^SP_r^S\cdots P_s^Sf-\Lambda^Sf\right|
 \le C\left\{M^{-7/24}\log M+{1\over M}
                   \sum_{i=r}^s(q_i+1)\right\}.            \tag{TP2}\label{eq:tu-TP2}
\end{equation}
The bound is uniform over the fixed realized radial word.  An arbitrary
inter-block reset terminates the estimate; the next block restarts from an
arbitrary law.
\end{lemma}

\begin{proof}[Proof of Lemma~\ref{lem:typed-centered-endpoint-product}]
The comparison is applied after averaging the exit edge, which removes the
roughness of the digital boundary. We then normalize each complete return
word and compare the resulting product with rotation-invariant Brownian
kernels.

\smallskip\noindent
\underline{\textit{Step 1.}}\quad Smooth the endpoint observables.
A discrete pair is
\((a,e)\in A\times E\).  Every exterior kernel starts at the head
\(\operatorname{out}(e)\) of the preceding exit edge.  With
\(P_i^S=J_i^SH^S\), exact operator composition is
\begin{equation}
 \Lambda^SP_r^S\cdots P_s^Sf
 =\lambda_A^SH^SJ_r^SH^S\cdots J_s^S({\cal H}^Sf).         \tag{TP3}\label{eq:tu-TP3}
\end{equation}
Thus the first operation on the terminal observable is the exact average
\({\cal H}^S\).

The continuum entrance and exit state spaces are separate copies of
\(\mathbb T\), in the common angular coordinate of the centered block.
Let \(\bar H\) be Brownian exit from radii \(M\) to \(eM\), and let
\(\bar J_i\) be the normalized Brownian return word.  Their densities are
rotation convolutions.  Explicitly,
\begin{equation}
 {d\bar H(\alpha,\cdot)\over d\theta/(2\pi)}(\beta)
 ={e^2-1\over e^2+1-2e\cos(\beta-\alpha)},                \tag{TP4}\label{eq:tu-TP4}
\end{equation}
and an elementary side kernel from logarithmic level \(x\) to \(b\) before
\(a\) has unnormalized density
\begin{equation}
 {1\over2\pi}\left\{{x-a\over b-a}
 +2\sum_{n\ge1}{\sinh(n(x-a))\over\sinh(n(b-a))}
       \cos n(\beta-\alpha)\right\}.                      \tag{TP5}\label{eq:tu-TP5}
\end{equation}
Consequently Haar measure \(m(d\alpha)=d\alpha/(2\pi)\) is preserved
by \(\bar H\) and every \(\bar J_i\). The continuum pair transition
\(\bar J_i\bar H\) therefore preserves
\(\bar\Lambda(d\alpha,d\beta)=m(d\alpha)\bar H(\alpha,d\beta)\).

We next construct the continuum test functions.  Write \(G^S=G_D\).
For \(x,x'\in D(y,M/2)\), the exact bridge identities are
\begin{equation}
 {\cal H}^Sm_x(a)=G^S(a,x),                               \tag{TP6}\label{eq:tu-TP6}
\end{equation}
\begin{equation}
 \begin{split}
 {\cal H}^Sq_{x,x'}(a)
 &=G^S(a,x)G^S(x,x')+G^S(a,x')G^S(x',x)\\
 &\quad-\mathbf1_{\{x=x'\}}G^S(a,x).
 \end{split}                                               \tag{TP7}\label{eq:tu-TP7}
\end{equation}
Let \(a_\alpha=y+Me^{i\alpha}\) and
\[
       \mathscr G^B(a_\alpha,x)={2\over\pi}
              g_{B(y,eM)}^B(a_\alpha,x).
\]
For \(f=m_x/C\), define
\(\varphi_f(\alpha)=\mathscr G^B(a_\alpha,x)/C\).  For
\(f=q_{x,x'}/[C(1+G^S(x,x'))]\), define \(\varphi_f\) by the right side
of \eqref{eq:tu-TP7}, divided by \(C(1+G^S(x,x'))\), and replace only the boundary-variable factors
\(G^S(a,x),G^S(a,x')\) by \(\mathscr G^B(a_\alpha,x)\) and
\(\mathscr G^B(a_\alpha,x')\).  The short-distance lattice coefficients
\(G^S(x,x')\) are retained, including on the diagonal.

After scaling by \(M\),
\begin{equation}
 g_{B(0,e)}^B(z,\xi)=\log\left|{e^2-z\bar\xi\over e(z-\xi)}\right|,
       \qquad |z|=1,\quad|\xi|\le\tfrac12,               \tag{TP8}\label{eq:tu-TP8}
\end{equation}
where the logarithmic Green function is normalized as in the preceding
display. The denominators stay uniformly
away from zero, and symmetry of \(G^S\) bounds all retained normalized
coefficients.  Hence
\begin{equation}
                  \sup_f\|\varphi_f\|_{C^4(\mathbb T)}\le C.          \tag{TP9}\label{eq:tu-TP9}
\end{equation}

We use the Green-function comparison of Kozdron--Lawler
\cite[Corollary~3.5, equation~(30)]{KozdronLawler2005}.
For a simply connected lattice domain \(A\) of inradius in \([n,2n]\),
let \(\widetilde A\) be its union-of-squares domain. If
\(u\in A_n^*:=\{z:g_{\widetilde A}(0,z)\ge n^{-1/16}\}\), the
comparison, away from the diagonal, is
\begin{equation}
 G_A(u,z)={2\over\pi}g_{\widetilde A}(u,z)+k_{z-u}
                 +O(n^{-7/24}\log n),                    \tag{TP10}\label{eq:tu-TP10}
\end{equation}
where \(a\) is the lattice potential kernel and
\(k_w=k_0+(2/\pi)\log|w|-a(w)\) for \(w\ne0\).
Apply this with the core pole in \(D(y,M/2)\). This pole is a fixed
fraction of the radius from the boundary and belongs to \(A_n^*\).
The boundary-variable point is at distance at least \(M/2-O(1)\) from
the pole, so \(k_{a-x}=O(M^{-2})\). Sandwiching the union-of-squares
disk between disks whose radii differ by \(O(1)\), and projecting
radially to the continuum circle, each costs \(O(M^{-1})\).
Equations
\eqref{eq:tu-TP6}--\eqref{eq:tu-TP10} therefore give, with
\(U_AF(a)=F(\arg(a-y))\),
\begin{equation}
 \|{\cal H}^Sf-U_A\varphi_f\|_\infty
       \le C\epsilon_G(M),\qquad
 \epsilon_G(M)=M^{-7/24}\log M.                           \tag{TP11}\label{eq:tu-TP11}
\end{equation}

For an exit edge put \(U_EF(e)=F(\arg(\operatorname{out}(e)-y))\).
Lemma~\ref{lem:smooth-interface-discretization}, with the same tubular
continuation for a disk having only its exit boundary, gives for every
smooth angular test \(F\)
\begin{equation}
 \|H^SU_EF-U_A\bar HF\|_\infty
       +|\lambda_A^SU_AF-mF|
       \le {C\over M}\|F\|_{C^4}.                         \tag{TP12}\label{eq:tu-TP12}
\end{equation}
For the second term, solve the disk Dirichlet problem with boundary data
\(F\), start at its center, and apply the tubular continuation argument
of (39e)--(39f). For the first term, solve the exit-disk problem and
start on the inner circle. This proves the comparison for the digital
circles directly.

\smallskip\noindent
\underline{\textit{Step 2.}}\quad Normalize each complete return word.
A prescribed word may have very small probability. The elementary
relative estimates must therefore be multiplied before the complete word
is normalized. Write its elementary kernels as
\[
 T_{i,l}^S(z,dz')=p_{i,l}^S(z)K_{i,l}^S(z,dz'),\qquad
 \bar T_{i,l}=\bar p_{i,l}\bar K_{i,l},\qquad1\le l\le q_i,
\]
and let \(Q_i^S=K_{i,1}^S\cdots K_{i,q_i}^S\).  Under this normalized
digital chain put
\[
 W_i=\prod_{l=1}^{q_i}{p_{i,l}^S(Z_{l-1})\over\bar p_{i,l}}.
\]
By (39c), if \(\eta_i=C\sum_lR_{i,l}^{-1}\le Cq_i/M\), then
\(e^{-\eta_i}\le W_i\le e^{\eta_i}\).  With
\(\bar s_i=\prod_l\bar p_{i,l}\), the identities
\begin{equation}
 {R_i^S1(e)\over\bar s_i}=E_{Q_{i,e}^S}W_i,\qquad
 J_i^SF(e)={E_{Q_{i,e}^S}[W_iF(Z_{q_i})]
                 \over E_{Q_{i,e}^S}W_i}                  \tag{TP13}\label{eq:tu-TP13}
\end{equation}
are exact.  Hence
\begin{equation}
             \|J_i^SF-Q_i^SF\|_\infty
             \le C\eta_i\|F\|_\infty.                    \tag{TP14}\label{eq:tu-TP14}
\end{equation}

Apply the normalized smooth-test estimate following (39c) to the
\(K_{i,l}^S\) and telescope from the right.  If
\(\Phi_{l,F}=\bar K_{i,l+1}\cdots\bar K_{i,q_i}F\), convolution gives
\(\|\Phi_{l,F}\|_{C^4}\le\|F\|_{C^4}\).  With the appropriate angular
pullback at every state,
\[
 \|Q_i^SU_AF-U_E\bar J_iF\|_\infty
       \le C\eta_i\|F\|_{C^4}.
\]
Together with \eqref{eq:tu-TP14},
\begin{equation}
 \|J_i^SU_AF-U_E\bar J_iF\|_\infty
       \le C\eta_i\|F\|_{C^4}.                            \tag{TP15}\label{eq:tu-TP15}
\end{equation}
The weight in \eqref{eq:tu-TP13} contains the row masses of every letter of the fixed
word, including letters after the state currently being compared.

\smallskip\noindent
\underline{\textit{Step 3.}}\quad Compare the complete product.

Write \eqref{eq:tu-TP3} as alternating operators
\(H^S,J_r^S,H^S,\ldots,J_s^S\), with terminal test
\(g_f^S={\cal H}^Sf\).  Let \(L_j^S,\bar L_j\) be the corresponding
discrete and Brownian operators. Write \(U_j\) for the angular pullback
on the state space after the \(j\)-th operator, and put
\(\Phi_{j,f}=\bar L_{j+1}\cdots\bar L_N\varphi_f\).  The operator identity, with every pullback on its corresponding state
space, is
\begin{equation}
\begin{split}
 L_1^S\cdots L_N^Sg_f^S-U_0\bar L_1\cdots\bar L_N\varphi_f
 &=L_1^S\cdots L_N^S(g_f^S-U_N\varphi_f)\\
 &\quad+\sum_{j=1}^NL_1^S\cdots L_{j-1}^S
 [L_j^SU_j\Phi_{j,f}-U_{j-1}\bar L_j\Phi_{j,f}].          \end{split}
\tag{TP16}\label{eq:tu-TP16}
\end{equation}
All \(\Phi_{j,f}\) have uniformly bounded \(C^4\)-norm by \eqref{eq:tu-TP9} and
convolution contraction.  Equations \eqref{eq:tu-TP11}, \eqref{eq:tu-TP12}, and \eqref{eq:tu-TP15} bound
\eqref{eq:tu-TP16} by the right side of \eqref{eq:tu-TP2}.  In the continuum the alternating
product preserves Haar flux, so it has integral \(m\varphi_f\).  A
separate zero-length use of \eqref{eq:tu-TP11}--\eqref{eq:tu-TP12} gives
\(|\Lambda^Sf-m\varphi_f|\le C(\epsilon_G+M^{-1})\).
Subtracting these two comparisons proves \eqref{eq:tu-TP2}.

For a lower conformal block the same proof uses its one common
\(\chi_y\)-angle.  Each occurrence of \(H^S\) has angular conjugacy error
\(O(M^{-1}+e^{-J+h})\).  Since \(2h\le J\),
\(e^{-J+h}\le M^{-1}\), and paying this error once per occurrence is
already contained in \(M^{-1}\sum_i(q_i+1)\). A decision depending on
an angular endpoint can change Haar measure. Such decisions are therefore
included in the reset density constructed below, at the boundary between
centered blocks.
\end{proof}

\subsection{Occupation of completed terminal excursions}
Fix a center \(y\) and set \(o_y:=y\).  Let
\[
\begin{aligned}
 D_y&=\{z\in\mathbb Z^2:\|z-y\|<eM\},\\
 A_y&=\partial_{\rm in}\{z:\|z-y\|<M\},\\
 V_y&=\{z:\|z-y\|<M/2\},
\end{aligned}
\]
where \(\partial_{\rm in}D=\{a\in D:\exists a'\notin D,\ a'\sim a\}\),
and let \(E_y=\{(u,v):u\in D_y,v\notin D_y,u\sim v\}\) be the directed
exit edges.  A completed
terminal piece begins at the first lattice vertex \(a\) after an inward
crossing of radius \(M\), and is killed on the directed edge \(e\in E_y\)
that first crosses radius \(eM\).  Erasing the interiors of all such
outward pieces gives exactly the intrinsic sigma-field \({\cal R}_y\) of
Lemma~\ref{lem:buffered-transfer}.  Conditional on their endpoint pairs
\(\omega_i=(a_i,e_i)\), their occupation vectors
\(X_i=(X_i(x):x\in V_y)\) are independent killed lattice bridges.

An \emph{admissible centered block word} \(w\) is a positive-probability
radial-only cylinder.  It has \(b=O(J^2)\) completed pieces, at most
\(Q\le J^3\) adjacent-shell decisions, and a fixed
number \(G_0\) of blocks.  Within each block every decision is between
circles centered at \(y\), or, in the lower spherical transfer, between successive level
sets of the single conformal coordinate \(\chi_y\).  In either case the
normalized Brownian kernels are rotation convolutions in one common
angular coordinate; this is what ``centered block'' means here.  Between
two lower blocks the safe connector contributes no observable and its
normalized angular kernel is treated only as an arbitrary reset.  Its
scalar prescribed-side row mass is retained in the exact word
factorization, but no comparison of its angular density or invariant law
is made.  Before factorization, retain only source intervals containing
at least one completed terminal piece as blocks.  Fold any initial empty
intervals and connectors into the first chronological law, every
intervening string of empty intervals and connectors into the preceding
reset, and every terminal such string into the final reset.  Hence
\(b_g\ge1\), and in the lower application \(G_0\le m_0\); the
inter-source connectors and empty intervals contribute no core
occupation.  In the upper application \(G_0=1\) when \(b>0\), after the
centered extension (39k).  The case \(b=0\) is trivial and is henceforth
excluded.  The reset construction below assigns every initial/final
suffix and connector exactly once, so there is no unrecorded future
weight.  Put
\[
 H(a,e)=P_a\big((S_{\tau_{D_y^c}-1},S_{\tau_{D_y^c}})=e\big).
\]
Write \(b_g\) for the number of completed pieces in block \(g\), so
\(\sum_gb_g=b\).  For \(i<b_g\), let
\(R_{g,i}^w(e,da')\) be the sub-Markov kernel of the prescribed exterior
return word, started at \(\operatorname{out}(e_{g,i})\), from
\(e_{g,i}=e\) to the next entrance
\(a_{g,i+1}=a'\).  The pair subkernel factors exactly as
\begin{equation}
 Q_{g,i}^w((a,e),d(a',e'))=
       R_{g,i}^w(e,da')H(a',de').                           \tag{47c}\label{eq:tu-47c}
\end{equation}
Write
\begin{equation}
 \sigma_{g,i}(e)=R_{g,i}^w1(e),\qquad
 J_{g,i}^w={R_{g,i}^w\over\sigma_{g,i}},
 \qquad\widetilde K_{g,i}^w=J_{g,i}^wH.                    \tag{47d}\label{eq:tu-47d}
\end{equation}
For \(g<G_0\), let
\begin{equation}
 C_g^w(e,d\zeta,d\omega')                                  \tag{47d'}\label{eq:tu-47d-prime}
\end{equation}
be the genuine sub-Markov reset kernel which, after the last endpoint
of block \(g\), contains the prescribed final no-more-piece suffix, the
safe connector and its data \(\zeta\), and the prefix of block \(g+1\)
through its first endpoint pair \(\omega'\).  For \(g=G_0\), the same
notation denotes the final suffix to a cemetery state.  Put
\begin{equation}
 \sigma_{g,b_g}(e)=C_g^w1(e),\qquad
 J_{g,b_g}^w=C_g^w/\sigma_{g,b_g}.                          \tag{47d''}\label{eq:tu-47d-doubleprime}
\end{equation}
In a centered suffix every adjacent-shell letter is included in \(Q\).
A lower inter-source connector is not declared centered; its normalized
kernel \(J_{g,b_g}^w\) is used as an arbitrary reset.

Let \(\alpha_{1,w}\) be the actual subprobability law of the first endpoint
pair after imposing only the word prefix through that piece.  Admissibility
means \(\alpha_{1,w}1>0\); put
\(\mu_{1,w}=\alpha_{1,w}/(\alpha_{1,w}1)\).  The
\emph{sequentially normalized} law \(\widetilde P_w\) is the
joint law of all endpoint and reset variables that starts with
\(\mu_{1,w}\), uses \(\widetilde K_{g,i}^w\) internally, and uses the
normalized reset kernels \(J_{g,b_g}^w\) between blocks and at the final
cemetery transition.  Thus the first law of every later block is
generated chronologically by the preceding reset, not conditioned on
its future output.  Write \(P_w\) for the true endpoint-and-reset law
conditioned on the complete word \(w\).  All estimates are uniform over
the arbitrary normalized reset kernels, which are marginalized at the
end.

Let \(\lambda_A\) be the law of the first hit of \(A_y\) by a walk started
at \(o_y\), and set
\begin{equation}
                 \Lambda_M(da,de)=\lambda_A(da)H(a,de).     \tag{47e}\label{eq:tu-47e}
\end{equation}
For a pair observable \(f\), write
\begin{equation}
                 ({\cal H}f)(a)=\sum_eH(a,e)f(a,e).         \tag{47f}\label{eq:tu-47f}
\end{equation}
For \(x,x'\in V_y\), define
\begin{equation}
 m_x(a,e)={G(a,x)H(x,e)\over H(a,e)},                       \tag{47g}\label{eq:tu-47g}
\end{equation}
\begin{equation}
\begin{aligned}
q_{x,x'}(a,e)=\frac{1}{H(a,e)}\big(&G(a,x)G(x,x')H(x',e)\\
&+G(a,x')G(x',x)H(x,e)\\
&-{\bf1}_{\{x=x'\}}G(a,x)H(x,e)\big).
\end{aligned}                                             \tag{47h}\label{eq:tu-47h}
\end{equation}
Uniformly,
\begin{equation}
 0\le m_x\le C,\qquad0\le q_{x,x'}\le C(1+G(x,x')),       \tag{47h'}\label{eq:tu-47h-prime}
\end{equation}
and the exit-edge averages are exactly
\[
 {\cal H}m_x(a)=G(a,x),
\]
\begin{equation}
 {\cal H}q_{x,x'}(a)=G(a,x)G(x,x')+G(a,x')G(x',x)
             -{\bf1}_{\{x=x'\}}G(a,x).                    \tag{47h''}\label{eq:tu-47h-doubleprime}
\end{equation}
The finite normalized observable family is
\begin{equation}
 {\cal F}_M=\{m_x/C:x\in V_y\}\cup
 \left\{{q_{x,x'}\over C(1+G(x,x'))}:x,x'\in V_y\right\}; \tag{47i}\label{eq:tu-47i}
\end{equation}
it has at most \(CM^4\) members.  Put
\begin{equation}
 \varepsilon_M=CM^{-7/24}\log M,
 \qquad
 \eta_w^*:=CQ/M+CG_{\rm nc},                                \tag{48a}\label{eq:tu-48a}
\end{equation}
where \(G_{\rm nc}\le G_0\le m_0\) counts noncentered whole resets; each
is factored once through the separator \eqref{eq:tu-59a0} below.  Thus
\(G_{\rm nc}=0\) in the upper application and is fixed in the lower one.
Take the deterministic centers \(c_f=\Lambda_Mf\).  Before the
normalizations in \eqref{eq:tu-47i}, they satisfy
\begin{equation}
 c_{m_x}=g+O(\varepsilon_M),                                \tag{48b}\label{eq:tu-48b}
\end{equation}
\begin{equation}
 c_{q_{x,x'}}=gG(x,x')+gG(x',x)-{\bf1}_{\{x=x'\}}g
       +O(\varepsilon_M(1+G(x,x'))).                        \tag{48c}\label{eq:tu-48c}
\end{equation}
\begin{lemma}\label{lem:terminal-feynman-kac}
With the preceding notation, the exact density of the complete word is
\begin{equation}
 {dP_w\over d\widetilde P_w}
 ={\displaystyle\prod_{g=1}^{G_0}\prod_{i=1}^{b_g}
          \sigma_{g,i}(e_{g,i})\over
   \displaystyle E_{\widetilde P_w}
       \prod_{g=1}^{G_0}\prod_{i=1}^{b_g}
          \sigma_{g,i}(e_{g,i})},
\qquad
 e^{-2\eta_w^*}\le
 {dP_w\over d\widetilde P_w}
 \le e^{2\eta_w^*}.                                        \tag{48d}\label{eq:tu-48d}
\end{equation}
The product visibly includes the final no-more-piece suffix and every
connector likelihood.  For every \(t>0\), conditionally on \(w\),
\begin{equation}
\begin{split}
P_w\Bigg(&\max_{f\in{\cal F}_M}\left|
 b^{-1}\sum_{i\le b}f(\omega_i)-c_f\right|\\
&>t+C\{G_0/b+\varepsilon_M+Q/M\}\Bigg)
 \le2|{\cal F}_M|\exp\{2\eta_w^*-cbt^2\}.                   \end{split}
\tag{49}\label{eq:tu-49}
\end{equation}

For \(x\in V_y\), \(G=G_{D_y}\), and
\begin{equation}
 r_\theta(x):=
 \frac{e^\theta-1}
 {1-(G(x,x)-1)(e^\theta-1)},                                 \tag{39}\label{eq:tu-39}
\end{equation}
uniformly for \(0\le r_\theta\le c/h\), \(h=\log M\),
\begin{equation}
 \log E\left[\exp\left\{\theta\sum_{i=1}^bX_i(x)\right\}
             \,\middle|\,w\right]
       =bg\,r_\theta+O\left(br_\theta^2+
       r_\theta\{G_0+b\varepsilon_M+bQ/M\}
                    +\eta_w^*\right).                       \tag{50}\label{eq:tu-50}
\end{equation}
The estimate is uniform in \(x,w,b\) and concerns completed pieces.
In both applications below, every visit to \(V_y\) lies in a completed
piece, so no conditional incomplete-piece estimate is needed.

For the upper-bound choice from Lemma~\ref{lem:buffered-transfer},
\[
 h=(C_h\log J)^4,\qquad n=J-h+O(1),\qquad
 U=u_{J,c},\quad c>2,
\]
put \(Y=\sqrt{2b}\).  On a coefficient-two coarse barrier, uniformly over
all \(0\le Y\le 3n\),
\begin{equation}
\begin{split}
 P\left(\max_{x\in V_y}\sum_iX_i(x)\ge U\mid w,b\right)
 &\le C e^{-2n}J^c e^{b/n}\\
 &\quad{}\times
 \exp\left\{-\kappa\min\left(
 \frac{(p-p_*)^2}{h},\frac{J^2}{h}\right)\right\},            \end{split}
\tag{51}\label{eq:tu-51}
\end{equation}
whenever \(Y\) is in the central range.  Here
\[
 p=B_n+1-Y,\qquad B_n=\rho_Jn+z_0+h^{1/4},\qquad
 p_*=O(\log J).
\]
Here the coefficient-two barrier means
\begin{equation}
 \sqrt{2T^S_{y,j}}\le
 B_j:=\rho_Jj+z_0+(j\wedge(J-j))^{1/4}
 \quad(k_y<j\le n),                                         \tag{52}\label{eq:tu-52}
\end{equation}
after the fixed \(O(1)\) relaxation of Lemma~\ref{lem:buffered-transfer}.  The central range is
\(\epsilon n\le Y\le B_n+1\), for one fixed small
\(\epsilon>0\).  For \(Y<\epsilon n\), \eqref{eq:tu-51} holds with its last factor
replaced by \(\exp(-\kappa J^2/h)\).
\end{lemma}

\begin{proof}[Proof of Lemma~\ref{lem:terminal-feynman-kac}]
The bridge formula reduces the occupation moment to a product of bounded
endpoint observables. We control that product under the sequential law,
where centered transitions contract, and then restore the likelihood of
the complete word. The likelihood comparison must include the last
suffix and every reset, because conditioning on their future outcomes
can otherwise change the endpoint distribution.

\smallskip\noindent
\underline{\textit{Step 1.}}\quad Compute bridge moments and their centers.
Let
\[
 H(a,e)=P_a\big((S_{\tau-1},S_\tau)=e\big),\qquad
 \tau=\tau_{D_y^c}.
\]
For the bridge from \(a\) to \(e\), write \(X(x)\) for its occupation of
\(x\).  Splitting at the first visit to \(x\), and then summing the
geometric number of returns to \(x\), gives the exact rank-one identity
\begin{equation}
\begin{split}
 H_\theta(a,e)
 &:=E_a[e^{\theta X(x)};(S_{\tau-1},S_\tau)=e]\\
 &=H(a,e)+r_\theta G(a,x)H(x,e).                              \end{split}
\tag{42}\label{eq:tu-42}
\end{equation}
Thus
\begin{equation}
       \frac{H_\theta(a,e)}{H(a,e)}
       =1+r_\theta m_x(a,e),\qquad
       m_x(a,e):=\frac{G(a,x)H(x,e)}{H(a,e)}.                 \tag{43}\label{eq:tu-43}
\end{equation}
A Harnack chain inside the fixed-ratio disk gives \(0\le m_x\le C\),
uniformly in \(x\in V_y\). Although \(m_x(a,e)\) is indexed by the
rough digital exit edges, the transition \eqref{eq:tu-47c} first averages over that
edge. From \eqref{eq:tu-43},
\begin{equation}
 ({\cal H}m_x)(a)=\sum_eH(a,e)m_x(a,e)=G(a,x).               \tag{43a}\label{eq:tu-43a}
\end{equation}
More generally, if \(q_{x,x'}(a,e)=E^{a,e}[X(x)X(x')]\), the
two ordered visit-time decompositions give
\begin{equation}
\begin{split}
 ({\cal H}q_{x,x'})(a)
 &=G(a,x)G(x,x')+G(a,x')G(x',x)\\
 &\quad-{\bf1}_{\{x=x'\}}G(a,x).                            \end{split}
\tag{43b}\label{eq:tu-43b}
\end{equation}
Thus the only functions subsequently fed into a return-word kernel are
Green functions with their singularities a distance at least \(M/2\)
from \(A_y\).

We next identify the common center and the only observables that require
comparison.  With \(\lambda_A\) defined as the first-hit law of the
inner vertex boundary, the strong Markov property gives
\begin{equation}
 d_x:=\sum_a\lambda_A(a)G_{D_y}(a,x)
 =G_{D(y,eM)}(o_y,x)-G_{D(y,M)}(o_y,x)
                  +O(M^{-1}\log M).                         \tag{44}\label{eq:tu-44}
\end{equation}
Indeed the one-edge collar can return from \(A_y\) to \(D(y,M/2)\)
before exiting \(D(y,M)\) with probability \(O(M^{-1})\), while the
Green function is at most \(C\log M\).  This remainder is
\(o(\varepsilon_M)\).

Apply the Green comparison \eqref{eq:tu-TP10} to both disks in \eqref{eq:tu-44}. For distinct
arguments it reads
\begin{equation}
 G_A(u,z)={2\over\pi}g_A(u,z)+k_{z-u}
                  +O(n^{-7/24}\log n),\tag{45}\label{eq:tu-45}
\end{equation}
with \(k_w=k_0+(2/\pi)\log|w|-a(w)\) for \(w\ne0\).
The lattice correction is identical in the two domains and cancels.
When \(x=o_y\), use instead the diagonal formula in
\cite[Theorem~1.2]{KozdronLawler2005}; its additive potential-kernel
constant also cancels. The core points lie a fixed fraction of the radius
from both boundaries, and the union-of-squares domains are between disks
whose radii differ by \(O(1)\). The Brownian Green-function difference
for the concentric disks is exactly \(\log e=1\). Consequently, uniformly
in \(x\in V_y\),
\begin{equation}
                     d_x=g+O(\varepsilon_M),
 \qquad \varepsilon_M=CM^{-7/24}\log M.                    \tag{46}\label{eq:tu-46}
\end{equation}
The identities \eqref{eq:tu-43a}--\eqref{eq:tu-43b} now give the centers \eqref{eq:tu-48b}--\eqref{eq:tu-48c}.

\smallskip\noindent
\underline{\textit{Step 2.}}\quad Compare and contract centered transitions.
The state of
\(R_{g,i}^w\) is the preceding directed exit edge, and the walk restarts at
its head; its target is the next inward directed edge.  A word is called
admissible only when its cylinder has positive probability, so every row
normalization below is defined.  Brownian row masses of its elementary
letters are scalar, while Lemma~\ref{lem:smooth-interface-discretization} gives
their relative lattice comparison.  The exact whole-word weighting formula
\eqref{eq:tu-TP13} therefore yields, for a scalar \(s_{g,i}>0\),
\begin{equation}
 \sup_e\left|{\sigma_{g,i}(e)\over s_{g,i}}-1\right|
       \le {Cq_i\over M}.                                  \tag{58c}\label{eq:tu-58c}
\end{equation}
This estimate holds for the complete word, since \eqref{eq:tu-TP13} retains the row
masses of all its letters.

In the common logarithmic conformal coordinate, every normalized Brownian
letter is the rotation convolution \eqref{eq:tu-TP5}, and the terminal disk exit is
\eqref{eq:tu-TP4}.  Apply Lemma~\ref{lem:typed-centered-endpoint-product} to the exact
exit-edge averages \eqref{eq:tu-43a}--\eqref{eq:tu-43b}.  For consecutive internal transitions
\(r\le s\) in one block and every \(f\in{\cal F}_M\), it gives
\begin{equation}
 \left|\Lambda_M
   \widetilde K_{g,r}^w\widetilde K_{g,r+1}^w\cdots
   \widetilde K_{g,s}^wf-c_f\right|
 \le C\left\{\varepsilon_M+
       {1\over M}\sum_{i=r}^s(q_i+1)\right\}.              \tag{58e}\label{eq:tu-58e}
\end{equation}
The two terms are compared separately with the same Haar integral of the
disk-Green traces \eqref{eq:tu-TP6}--\eqref{eq:tu-TP9}. The estimate therefore controls complete
prefixes even though \(\Lambda_M\) need not be invariant for a single
discrete transition. In a lower source interval the common coordinate is
\(\chi_y\); its terminal chart error is included in the displayed sum.
An inter-source connector ends this comparison.

To remove dependence on the initial law of each block, we also need a
uniform contraction in total variation. The elliptic Harnack inequality
\cite[Theorem~1.1]{AthreyaEtAl2022} states that for every \(d\ge1\)
there is \(C_d\) such that a nonnegative function harmonic in
\(B(x_0,R)\) satisfies
\begin{equation}
       \max_{B(x_0,R/2)}u\le C_d\min_{B(x_0,R/2)}u.           \tag{58a}\label{eq:tu-58a}
\end{equation}
Starting one step inside the outer boundary is too close to use \eqref{eq:tu-58a}
directly.  Insert instead the deterministic intermediate circle
\[
                 C_M=\partial D(y,e^{1/2}M),
\]
which every successful outer-to-inner crossing must hit.  For a target
atom \(a\in A_y\), put
\[
 h_a(z)=P_z(S_{H_{A_y}}=a,\ H_{A_y}<\tau_{D(y,eM)}),
 \qquad s(z)=\sum_{a\in A_y}h_a(z),\qquad z\in C_M.
\]
A bounded chain of lattice balls of radius \(cM\), all a fixed positive
fraction from both boundary components, applies \eqref{eq:tu-58a} to \(h_a\) and
\(s\).  Uniformly in \(a,z,z_0\in C_M\),
\[
 ch_a(z_0)\le h_a(z)\le Ch_a(z_0),\qquad
 cs(z_0)\le s(z)\le Cs(z_0).
\]
After division by the total successful-leg probability this gives
\begin{equation}
       J_{\rm in}(z,\cdot)\ge\alpha
                     J_{\rm in}(z_0,\cdot)                  \tag{58b}\label{eq:tu-58b}
\end{equation}
where \(J_{\rm in}(z,\cdot)=h_{\cdot}(z)/s(z)\) is the entrance
law conditional on reaching the inner boundary, and \(\alpha>0\) is
uniform in \(M,z,z_0\). Precomposition by all earlier
within-block word factors and postcomposition by \(H\) preserve this row-independent
minorization. Writing \(\operatorname{osc}(f)=\sup f-\inf f\), we obtain
\(\operatorname {osc}(\widetilde K_{g,i}^wf)\le(1-\alpha)
\operatorname {osc}(f)\) for every bounded \(f\) at every within-block
transition.  Write \(\rho=1-\alpha<1\).  A reset transition is used
exactly and has Dobrushin coefficient at most one; there are only
\(G_0-1\) of them.

\smallskip\noindent
\underline{\textit{Step 3.}}\quad Retain the full likelihood at resets.
A connector can have an arbitrary normalized angular kernel. What is
needed is a bound on the oscillation of its row mass, uniform even when
the prescribed future event is rare. To obtain it, factor each reset
through a deterministic separator. Fix
\begin{equation}
        \lambda_*=e^{3/2}\in(e,e^2),\qquad
        \Sigma_y=\partial D(o_y,\lambda_*M).                \tag{59a0}\label{eq:tu-59a0}
\end{equation}
The next active outer interface after the terminal \(eM\)-circle has scale
\(e^2M\).  The fixed-gap geometry and \eqref{eq:tu-47c} put \(\Sigma_y\) strictly
between them and keep \(D(o_y,\lambda_*M)\) inside the applicable source
interval and global killing domain.

Let \(e=(u,v)\in E_y\) be the last exit edge of a nonempty block, so the
exterior kernel starts at \(v=\operatorname{out}(e)\) and
\(|v-o_y|=eM+O(1)\).  Put
\[
 \theta_-=H_{D(o_y,M)},\qquad
 \theta_+=\tau_{D(o_y,\lambda_*M)},
\]
and on \(\{\theta_+<\theta_-\}\) let \(F_+\) be the directed outward
edge at \(\theta_+\).  For \(g<G_0\), let
\(B_g^w(f,d\zeta,d\omega')\) be the genuine post-separator subkernel
which realizes the rest of the no-more-piece suffix, all folded empty
source intervals and safe connector signs, and the next nonempty block
prefix through its first endpoint pair.  For \(g=G_0\) it instead contains
the remaining suffix and final cemetery transition.  Durations, parities,
angles, and already visible edge data are outputs of this kernel, never
extra scalar conditions.  Set \(r_g^w(f)=B_g^w(f,\hbox{all outputs})\).
The strong Markov property gives the exact factorization
\begin{equation}
 C_g^w(e,d\gamma)=\sum_f
 P_v(\theta_+<\theta_-,F_+=f)B_g^w(f,d\gamma).              \tag{59a1}\label{eq:tu-59a1}
\end{equation}
Indeed, a return to \(D(o_y,M)\) before \(\Sigma_y\) after the alleged
last completed piece would, on the subsequent trip to the block's outer
endpoint, force another exit of \(D_y\), hence another completed terminal
piece.  Conversely, no reset restriction other than avoiding the inner
circle is tested before \(\Sigma_y\); all remaining restrictions occur in
\(B_g^w\).

Define on the digital annulus
\[
 \Omega_y=D(o_y,\lambda_*M)\setminus D(o_y,M)
\]
\begin{equation}
 h_g^w(z)=E_z[r_g^w(F_+);\theta_+<\theta_-].                \tag{59a2}\label{eq:tu-59a2}
\end{equation}
This is a nonnegative discrete harmonic function at every interior vertex,
with exact Dirichlet value zero on the inner boundary and the possibly rough
edge data \(r_g^w\) on the outer boundary.  Global killing and the entire
later connector/prefix occur only in that outer boundary datum.  Moreover,
\begin{equation}
                  \sigma_{g,b_g}(e)=C_g^w1(e)=h_g^w(v).     \tag{59a3}\label{eq:tu-59a3}
\end{equation}
All such \(v\)'s lie in the band
\(\bigl|\|v-o_y\|-eM\bigr|\le2\), a distance
comparable to \(M\) from both components of \(\partial\Omega_y\).
A fixed number of overlapping lattice balls of radius \(cM\), whose doubled
balls remain in \(\Omega_y\), connects any two points of this band.
Applying \eqref{eq:tu-58a} along this chain gives
\begin{equation}
       \sup_e\sigma_{g,b_g}(e)
              \le C_H\inf_e\sigma_{g,b_g}(e).              \tag{59a}\label{eq:tu-59a}
\end{equation}
If \(h_g^w\) vanishes at one interior start, the mean-value property and
connectedness make it identically zero, and the complete word cylinder has
zero mass and is discarded.  Otherwise every row is positive and the
normalization \(J_{g,b_g}^w=C_g^w/\sigma_{g,b_g}\) is defined.  This proves
\eqref{eq:tu-59a} independently of how rare or rough the later event is.

Start the first block with its actual chronological law through the first
endpoint.  Use every normalized internal kernel and every normalized reset
in chronological order, including the final cemetery output, to define
\(\widetilde P_w\).  Direct disintegration of the unnormalized cylinder
measure gives
\begin{equation}
 {dP_w\over d\widetilde P_w}
 ={\displaystyle\prod_{g=1}^{G_0}\prod_{i=1}^{b_g}
          \sigma_{g,i}(e_{g,i})\over
   \displaystyle E_{\widetilde P_w}
       \prod_{g=1}^{G_0}\prod_{i=1}^{b_g}
          \sigma_{g,i}(e_{g,i})}.                           \tag{60a}\label{eq:tu-60a}
\end{equation}
The initial prefix is absorbed into \(\mu_{1,w}\), and the product
includes the final suffix. Equation \eqref{eq:tu-58c}, applied also to a centered final suffix, bounds the
logarithmic row oscillation of centered factors by \(Cq_i/M\); \eqref{eq:tu-59a} costs
\(\log C_H\) once per genuinely noncentered reset.  Thus in the upper
centered extension \(G_{\rm nc}=0\), while in the lower construction
\(G_{\rm nc}\le G_0\le m_0\).  Scalar
constants cancel between numerator and denominator, so, after absorbing
\(\log C_H\) into \(C\),
\begin{equation}
 \left|\log{dP_w\over d\widetilde P_w}\right|
 \le2\{CQ/M+CG_{\rm nc}\}=2\eta_w^*.                       \tag{60b}\label{eq:tu-60b}
\end{equation}
\smallskip\noindent
\underline{\textit{Step 4.}}\quad Concentrate the endpoint averages.
Fix \(w\) and a normalized
\(f\in{\cal F}_M\), expose the sequential joint chain under
\(\widetilde P_w\), and use the Doob martingale of
the whole sum \(\sum_{i\le b}f(\omega_i)\).  Changing \(\omega_i\) changes the
conditional expectation of the future sum by at most
\[
       2G_0\sum_{j\ge0}\rho^j\le {2G_0\over1-\rho};
\]
the factor \(G_0\) accounts for the possibility that a reset preserves
all the remaining variation.
Azuma--Hoeffding consequently gives
\begin{equation}
 \widetilde P_w\left(\left|\sum_{i\le b}f(\omega_i)
       -E_{\widetilde P_w}\sum_{i\le b}f(\omega_i)\right|>bt\right)
 \le2e^{-cbt^2}.                                             \tag{61}\label{eq:tu-61}
\end{equation}
More explicitly, if the \(k\)-th endpoint of a block has actual initial
law \(\mu\), contraction and the full prefix comparison \eqref{eq:tu-58e} give
\begin{equation}
\begin{split}
\left|\mu\widetilde K_{g,1}^w\cdots
 \widetilde K_{g,k-1}^wf-c_f\right|
 &\le2\rho^{k-1}\\
 &\quad+C\left\{\varepsilon_M+
 {1\over M}\sum_{j<k}(q_j+1)\right\}.                      \end{split}
\tag{61a0}\label{eq:tu-61a0}
\end{equation}
For \(k=1\) this just uses \(\|f\|_\infty\le1\).
Summing \eqref{eq:tu-61a0} over all endpoints and blocks, and using
\(\sum_j(q_j+1)\le C Q\), gives
\begin{equation}
 \left|E_{\widetilde P_w}\sum_{i\le b}f(\omega_i)-bc_f\right|
 \le C\{G_0+b\varepsilon_M+bQ/M\}.                         \tag{61a}\label{eq:tu-61a}
\end{equation}
A union bound under \(\widetilde P_w\), followed by \eqref{eq:tu-48d}, proves
\eqref{eq:tu-49}.

\smallskip\noindent
\underline{\textit{Step 5.}}\quad Bound the exponential occupation moment.
Conditional on the endpoints, \eqref{eq:tu-42}--\eqref{eq:tu-43} makes the exponential moment
\(\prod_i(1+rm_x(\omega_i))\).  Under \(\widetilde P_w\), put
\begin{equation}
                 \psi_i=\log(1+rm_x(\omega_i)).              \tag{62}\label{eq:tu-62}
\end{equation}
Since \(0\le m_x\le C\),
\(\psi_i=rm_x(\omega_i)+O(r^2)\) and
\(\operatorname {osc}(\psi_i)\le Cr\).  The same influence calculation
as for \eqref{eq:tu-61} shows that every Doob martingale difference for
\(\sum_i\psi_i\) is at most \(CG_0r/(1-\rho)\).  Conditional Hoeffding and
Jensen therefore give
\begin{equation}
 0\le\log E_{\widetilde P_w}e^{\sum_i\psi_i}
       -E_{\widetilde P_w}\sum_i\psi_i\le Cbr^2.             \tag{62a}\label{eq:tu-62a}
\end{equation}
The exact exit-edge average and its Taylor remainder are
\begin{equation}
 \sum_eH(a,e)\log(1+rm_x(a,e))=rG(a,x)+O(r^2),             \tag{62b}\label{eq:tu-62b}
\end{equation}
because \(m_x\le C\) and \(\sum_eH(a,e)m_x(a,e)=G(a,x)\le C\).
The product telescoping \eqref{eq:tu-61a}, now for \(m_x/C\), and \eqref{eq:tu-46} give
\begin{equation}
 \sum_{i\le b}E_{\widetilde P_w}m_x(\omega_i)
       =bg+O\{G_0+b\varepsilon_M+bQ/M\}.                    \tag{62c}\label{eq:tu-62c}
\end{equation}
Combining \eqref{eq:tu-62}--\eqref{eq:tu-62c} proves the completed-piece version of \eqref{eq:tu-50} under
\(\widetilde P_w\).  Finally,
\eqref{eq:tu-48d} puts the exact full-word density \eqref{eq:tu-60a} between
\(e^{-2\eta_w^*}\) and \(e^{2\eta_w^*}\).  Multiplying the nonnegative
exponential by this global density and taking logarithms changes the
answer by at most \(2\eta_w^*\), exactly the last error in \eqref{eq:tu-50}.
This proves \eqref{eq:tu-50} under \(P_w\).

\smallskip\noindent
\underline{\textit{Step 6.}}\quad Optimize the terminal upper tail.
The first choice of tilt will cancel the endpoint Gaussian cost. A further
adjustment supplies decay in the distance from the preferred endpoint,
which is needed when summing all endpoint bands. The inverse of \eqref{eq:tu-39} is
\begin{equation}
       \theta(r)=
       \log\frac{1+G(x,x)r}{1+(G(x,x)-1)r}.                   \tag{49}\label{eq:tu-49-2}
\end{equation}
Take \(r_0=(gn)^{-1}\).  Since
\(G(x,x)=gh+O(1)\) and \(h^2/J=o(1)\),
\begin{equation}
 \theta(r_0)=\frac1{gJ}+O(J^{-2}+h^2J^{-3}),\qquad
 \theta(r_0)U=2J-c\log J+O(1),\qquad bg r_0=b/n.              \tag{50}\label{eq:tu-50-2}
\end{equation}
Since \(G_0=1\), \(b=O(J^2)\), \(Q\le J^3\), and
\(M=e^{(C_h\log J)^4}\), every error term in \eqref{eq:tu-50}
is \(O(1)\) at \(r_0\), uniformly in the word.
A union bound over \(O(e^{2h})\) sites, \eqref{eq:tu-50}, and Chernoff's inequality
give the first line of \eqref{eq:tu-51}.

We retain the tilt freedom to sum every endpoint band.  Set
\[
       F_b(r)=\theta(r)U-bgr-C_0br^2.
\]
Here are the derivative and feasibility details.  Uniformly on every fixed
interval \(0\le r\le a/h\),
\[
 \theta'(r)=\frac1{(1+Gr)(1+(G-1)r)},
\]
\[
 -\theta''(r)=\theta'(r)
 \left\{\frac G{1+Gr}+\frac{G-1}{1+(G-1)r}\right\}.
\]
Thus \(c_a\le\theta'(r)\le C_a\),
\(c_ah\le-\theta''(r)\le C_ah\), and
\(|\theta'''(r)|\le C_ah^2\).  Since \(b=O(J^2)\), after increasing
\(J\),
\begin{equation}
       -C_1J^2h\le F_b''(r)\le-c_1J^2h.                     \tag{51a}\label{eq:tu-51a}
\end{equation}
Also
\begin{equation}
 \theta'(r_0)=1-\frac{2h}{n}
                  +O(n^{-1}+h^2n^{-2}).                     \tag{51b}\label{eq:tu-51b}
\end{equation}
Define \(b_*\) by \(F_{b_*}'(r_0)=0\).  Since
\[
       b_*={U\theta'(r_0)\over g+2C_0r_0},
\]
the relations \(J=n+h+O(1)\) and \(h^2=o(J)\) give
\[
 b_*=2n^2-cn\log J+O(n+h^2+h\log J)
     =2n^2-cn\log J+O(n).
\]
Consequently
\begin{equation}
       Y_*=\sqrt{2b_*}
          =2n-\frac c2\log J+O(1),                            \tag{51}\label{eq:tu-51-2}
\end{equation}
and hence \(p_*=B_n+1-Y_*=O(\log J)\).  More exactly,
\begin{equation}
 F_b'(r_0)=(g+2C_0r_0)(b_*-b)
          =(g+o(1)){(Y_*-Y)(Y_*+Y)\over2}.                  \tag{51c}\label{eq:tu-51c}
\end{equation}

Fix a sufficiently small constant \(\delta>0\).  If
\(\epsilon n\le Y\le B_n+1\) and \(|Y-Y_*|\le\delta J\), Taylor's
formula, \eqref{eq:tu-51a}, and the displacement
\[
       s={F_b'(r_0)\over C_1J^2h}
\]
(with \(C_1\) enlarged once) give
\begin{equation}
 \sup_rF_b(r)-F_b(r_0)
 \ge\kappa\frac{(Y-Y_*)^2}{h}.                              \tag{51d}\label{eq:tu-51d}
\end{equation}
This displacement is feasible.  If it is negative, the barrier and
\(h^{1/4}=O(\log J)\) imply \(Y-Y_*=O(\log J)\), whence
\(|s|=O(\log J/(Jh))=o(r_0)\).  If it is positive, choosing \(\delta\)
small ensures \(s<a/(2h)\).  When
\(\epsilon n\le Y\le Y_*-\delta J\), choose instead a fixed
\(s=a_0/h>0\), with \(a_0\) small compared with \(\delta/C_1\).
Equations \eqref{eq:tu-51a}--\eqref{eq:tu-51c} then give
\begin{equation}
       F_b(r_0+s)-F_b(r_0)\ge\kappa J^2/h.                  \tag{51e}\label{eq:tu-51e}
\end{equation}
There is no range \(Y\ge Y_*+\delta J\) because the upper barrier permits
only an \(O(\log J)\) overshoot above \(Y_*\).

Finally, if \(Y<\epsilon n\), take \(r=a_0/h\) directly.  The integral
formula \(\theta(r)=\int_0^r\theta'(t)\,dt\) gives
\(\theta(a_0/h)\ge c_aa_0/h\).  First choosing \(\epsilon\) and then
\(a_0\) small yields
\begin{equation}
       F_b(a_0/h)\ge\kappa_0J^2/h,
       \qquad F_b(r_0)=O(J)=o(J^2/h).                        \tag{51f}\label{eq:tu-51f}
\end{equation}
Equations \eqref{eq:tu-51d}--\eqref{eq:tu-51f}, and
\(p-p_*=Y_*-Y\), prove precisely the two alternatives in \eqref{eq:tu-51}.

\smallskip\noindent
\underline{\textit{Step 7.}}\quad Account for the initial and final portions.
In the applications, all core occupation belongs to completed pieces.  A lower source block begins and ends on interfaces outside
\(D_y\).  An upper block with \(k\le C_0\log J\) starts at the origin,
which is outside \(D_y\), and ends at \(\sigma_y\), also outside
\(D_y\).  Before the first inward crossing of radius \(M\) the path
cannot visit \(V_y\).  After any visit to \(V_y\), reaching the block's
outer endpoint forces a later crossing of radius \(eM\), so that inward
crossing and all intervening core occupation form a completed piece.
Thus the initial and final incomplete portions contribute exactly zero
to every \(X(x)\) used here.  Upper cells with \(k>C_0\log J\), including
the cells near the walk's starting point, are treated separately by the
one-site Green estimate in \eqref{eq:tu-62-2}.
\end{proof}

\subsection{A local upper bound for the maximum}
\begin{proposition}\label{prop:localized-disk-upper}
For every fixed \(\beta>0,c>2\) and \(A<\infty\),
\begin{equation}
\begin{split}
 P\left(\max_{\lVert x\rVert\le e^JJ^{-\beta}}
             L_{\tau_{e^J}}(x)\ge u_{J,c}\right)
 \le J^{-2\beta+c-2+o(1)}+O(J^{-A}).                         \end{split}
\tag{53}\label{eq:tu-53}
\end{equation}
The \(o(1)\) is uniform when \(\beta,c\) range over compact subsets of
\((0,\infty)\) and \((2,\infty)\), respectively.
\end{proposition}

\begin{proof}[Proof of Proposition~\ref{prop:localized-disk-upper}]
We work with one terminal cell at a time and extend the walk to a disk
centered at that cell. This puts the source law, traversal barrier, and
terminal word in the same centered geometry. The traversal estimate and
the terminal exponential moment then have opposite Gaussian factors;
after their cancellation, the ballot prefactor supplies the required
power \(J^{-2}\).

\smallskip\noindent
\underline{\textit{Step 1.}}\quad Choose cells and exclude barrier crossings.
Use the upper buffer
\[
       h=(C_h\log J)^4,\qquad M=e^h,\qquad n=J-h,
\]
with \(C_h\) fixed large.  Cover the target disk by original \(M\)-cells
obtained from the dilated image of Rosen's predetermined net \(F_n\),
rounding each image center to its nearest lattice point.  The chart has
uniformly bounded distortion on the relevant compact patch, so these
centers are \(cM\)-separated and form a \(CM\)-cover.  After altering the
fixed core and terminal radius factors, every lattice point lies in a
core, the overlap multiplicity is bounded, and Lemmas
\ref{lem:buffered-transfer} and \ref{lem:terminal-feynman-kac}
apply to the same cells.  We continue to denote their radii by \(M/2\) and \(eM\).
For a cell center with \(\|y\|>M\), use the exact index from (45a0),
\begin{equation}
 u_y=\log {3e^J\over\|y\|},\qquad
 k=k(y)=\lceil u_y\rceil,\qquad m=k+1.                     \tag{54a}\label{eq:tu-54a}
\end{equation}
Cells with \(\|y\|\le M\) are assigned to the final sitewise range
\(k>C_0\log J\).  Uniform distortion of the fixed chart makes \eqref{eq:tu-54a}
equal to the spherical source index (24) up to a bounded additive error,
which is absorbed by the barrier slack.
Every target cell has
\begin{equation}
       k\ge s:=\beta\log J+O(1).                              \tag{54}\label{eq:tu-54}
\end{equation}
For this fixed cell continue the walk from \(\tau_{e^J}\) to
\(\sigma_y=\tau_{D(y,3e^J)}\).  By (39l), a high local time in the cell
before \(\tau_{e^J}\) is also high before \(\sigma_y\).  All counts,
words, terminal bridges, and the source index below refer to this single
\(y\)-centered extended path.  This domination is applied separately to
each cell before the final union bound.

First exclude a crossing of the fixed positive-intercept
coefficient-two barrier \eqref{eq:tu-52} at or before level \(n\).  Since
\(h>(4\log J)^4\) after choosing \(C_h\) large, all these levels lie in
the range of (45d)--(45e).  Restricting the source shells to \(k\ge s\)
and using (40) gives
\begin{equation}
\begin{split}
 P(\hbox{some bad prefix in the target})
 &\le C\sum_{k\ge s}\sum_{l=k+1}^{n}
 e^{2(l-k)+\frac12 l_J^{1/4}}
 e^{-2l-l_J^{1/4}}\\
 &\le J^{-2\beta+o(1)},\qquad
 l_J:=l\wedge(J-l).                                          \end{split}
\tag{55}\label{eq:tu-55}
\end{equation}
The direct lattice estimates (45a)--(45e), for both the representative
count and the probability of a crossing at one center, apply through the
source-shell range \(k\le C_0\log J\).

\smallskip\noindent
\underline{\textit{Step 2.}}\quad Include the source law in the ballot estimate.
Fix one original cell in a shell
\[
       s\le k\le C_0\log J.
\]
Sum the complete words that include the source and whose radial prefix
stays below the barrier. The comparison is applied to this joint event
before conditioning on a particular source height.  Put \(m=k+1\) and
\[
       X=\sqrt{2T_m},\qquad Y=\sqrt{2T_n}.
\]
The directed-edge source estimate (45c) gives, for each unit band \(I_x\),
\begin{equation}
       P(X\in I_x)\le C(k+1)^C\exp\{-x^2/(2m)\}.             \tag{56}\label{eq:tu-56}
\end{equation}
For a unit endpoint band \(I_q=[q,q+1]\), let \(\mathscr A(q)\) denote
the disjoint union of complete words that retain the good barrier and end
in \(Y\in I_q\).  The joint source-inclusive comparison (45g), before any
disintegration of the actual lattice law, gives for every source band
\(I_x\)
\begin{equation}
\begin{split}
 &P_S\bigl(X\in I_x,\mathscr A(q)\bigr)\\
 &\qquad\le e^{CJ^3/M}
   \sum_{\sqrt{2t}\in I_x}\pi_{m,u_y}(t)
       P_t^{\rm GW}(\mathscr A_t^{\rm GW}(q)).              \end{split}
\tag{56a}\label{eq:tu-56a}
\end{equation}
Only the ideal law on the right is disintegrated by \(t\).  The remaining
traversal chain has length \(L=n-m\).  Its affine barrier has endpoints
\(a=B_m\) and \(b=B_n+1\), and (31a) shows that it dominates the truncated
\(J\)-barrier without replacing \(\rho_J\) by \(\rho_n\).

Fix \(\epsilon_0>0\) small and
\(\epsilon_0n\le q\le B_n+1\), and put \(p=B_n+1-q\).  For each
\(t\ge2\), set \(x_t=\sqrt{2t}\) and apply (30) pointwise with initial
height \(x_t\), endpoint band \([q,q+1]\), and duration \(L\).  On the
good prefix, \(x_t\le B_m=O(m)\), while the exact geometric source law
satisfies
\[
       \pi_{m,u_y}(t)\le Cm^{-1}\exp\{-x_t^2/(2m)\}.
\]
Because \(m+L=n\), the two exponential costs combine by the exact identity
\begin{equation}
 \frac{x_t^2}{m}+\frac{(q-x_t)^2}{L}
 =\frac{q^2}{n}+\frac{n(x_t-mq/n)^2}{mL}.                   \tag{57}\label{eq:tu-57}
\end{equation}
The identity uses the actual height \(x_t\). Grouping those heights into
unit bands changes the exponential by only \(e^{O(1)}\).  The Gaussian sum in the
last term and the remaining source factors cost only a polynomial in
\(m\).  Since \(q\ge\epsilon_0n\) and \(L\asymp n\), the ballot prefactor
in (30), including \(\sqrt{x_t/(qL)}\), gives
\begin{equation}
 P(\hbox{prefix good},\,Y\in I_q)
 \le C(k+1)^C(1+p)n^{-2}\exp\{-q^2/(2n)\}.                  \tag{58}\label{eq:tu-58}
\end{equation}
For the only positive start below two, namely \(T_m=1\), add an
independent ancestor and require that ancestor to die in its first
generation.  This has probability \(1/2\), so the desired probability is
at most twice the bound from initial population two: from generation
one onward the two paths agree, and these are exactly the generations on
which the barrier in (30) is imposed.  Initial population zero cannot
produce the present positive endpoint.  Thus
\eqref{eq:tu-58} proves the central-range estimate (41). Its relative lattice
comparison remains valid after summing over the cells.

\smallskip\noindent
\underline{\textit{Step 3.}}\quad Combine the traversal and terminal costs.
The complementary endpoint range is controlled directly by the terminal
tilt.  If
\(Y<\epsilon_0n\), then \(b\le\epsilon_0^2n^2/2+O(n)\), and the saturated
line of \eqref{eq:tu-51}, uniformly in the radial word, gives
\begin{equation}
 P(\hbox{terminal mark}\mid w)
 \le C e^{-2n}J^c
       \exp\{\epsilon_0^2n/2-\kappa J^2/h+O(1)\}.            \tag{58a}\label{eq:tu-58a-2}
\end{equation}
After summing this over at most \(Ce^{2n}\) original cells, the result is
\(o(J^{-A})\) for every fixed \(A\), because \(h=o(J)\), after fixing
\(\epsilon_0\) sufficiently small.  This is why neither (29) nor (41)
was claimed in the extinction range.

For the number \(b=T_n+O(1)\) of completed outward pieces in the band
\(I_q\),
\[
       b=q^2/2+O(n).
\]
Conditional on the radial word, Lemma~\ref{lem:terminal-feynman-kac} applies;
the start/end observation at the end of that lemma shows that all visits
to \(V_y\) belong to its completed pieces.  Multiply
\eqref{eq:tu-58} by \eqref{eq:tu-51}.  The factors
\(\exp\{-q^2/(2n)\}\) and \(\exp\{b/n\}\) cancel up to a constant.
Moreover,
\begin{equation}
 \sum_{0\le p\le B_n+1}(1+p)
 \exp\left\{-\kappa\min\left(
          \frac{(p-p_*)^2}{h},\frac{J^2}{h}\right)\right\}
       =J^{o(1)}.                                             \tag{59}\label{eq:tu-59}
\end{equation}
Indeed the nonsaturated Gaussian sum is
\[
 O\bigl((1+|p_*|)\sqrt h+h\bigr)=J^{o(1)},
\]
while its saturated complement contributes
\(O(J^2)e^{-\kappa J^2/h}=o(1)\).  The small-\(Y\) part has the
saturated factor in \eqref{eq:tu-51}.  Thus, uniformly
over all endpoint bands and all radial words,
\begin{equation}
 P(\hbox{this original cell is good and marked})
       \le (k+1)^Ce^{-2n}J^{c-2+o(1)}.                        \tag{60}\label{eq:tu-60}
\end{equation}
The terminal maximum has already included all sites of the original
cell, so the remaining union bound is over the original cells.

There are \(O(e^{2(n-k)})\) original cells in shell \(k\).  Summing \eqref{eq:tu-60}
and using \eqref{eq:tu-54} yields
\begin{equation}
 \sum_{k=s}^{C_0\log J}e^{2(n-k)}
       (k+1)^Ce^{-2n}J^{c-2+o(1)}
       \le J^{-2\beta+c-2+o(1)}.                              \tag{61}\label{eq:tu-61-2}
\end{equation}

\smallskip\noindent
\underline{\textit{Step 4.}}\quad Treat cells near the starting point.
It remains to treat \(k>C_0\log J\). For a lattice site in shell \(k\),
Green-function identities and Lemma~\ref{lem:elementary-walk-estimates} give
\[
 P_0(H_x<\sigma_y)\le C(k+1)/J
\]
and, conditionally on the hit,
\[
 P_x(L_{\sigma_y}(x)\ge u_{J,c})
       \le Ce^{-2J}J^c.
\]
There are \(O(e^{2J-2k})\) sites in that shell.  Their total contribution
is
\begin{equation}
       J^{c-1-2C_0+o(1)}.                                    \tag{62}\label{eq:tu-62-2}
\end{equation}
Choose \(C_0\) after \(A,c\).  Combining \eqref{eq:tu-55}, \eqref{eq:tu-61-2}, and \eqref{eq:tu-62-2} proves
\eqref{eq:tu-53}.  The word-cap and digital row errors are superpolynomial by
Lemma~\ref{lem:buffered-transfer} and are included in \(O(J^{-A})\).
\end{proof}
